% TikZ 图：质谱仪原理
% Day 22：洛伦兹力
\begin{tikzpicture}[scale=0.82,>=Latex,line cap=round,line join=round]
  \path[use as bounding box] (-5.10,-1.75) rectangle (3.45,3.55);

  % 速度选择器：E 与 B_1 正交
  \draw[thick] (-4.75,1.15) -- (-1.05,1.15);
  \draw[thick] (-4.75,-1.15) -- (-1.05,-1.15);
  \node at (-4.52,1.42) {$+$};
  \node at (-4.52,-1.42) {$-$};
  \foreach \x in {-4.15,-3.35,-2.55,-1.75} {
    \draw[->,blue,thick] (\x,0.88) -- (\x,-0.88);
  }
  \node[blue,font=\small] at (-1.42,0.84) {$E$ 向下};
  \foreach \x in {-3.82,-3.02,-2.22,-1.42} {
    \foreach \y in {-0.48,0.48} {
      \draw[gray,thick] (\x-0.055,\y-0.055) -- (\x+0.055,\y+0.055);
      \draw[gray,thick] (\x-0.055,\y+0.055) -- (\x+0.055,\y-0.055);
    }
  }
  \node[gray,font=\small] at (-4.18,0.76) {$B_1$ 向里};
  \node[font=\small] at (-2.90,-1.62) {速度选择器：$v=E/B_1$};

  \fill[red] (-4.65,0) circle (2.2pt);
  \node[red,above left] at (-4.65,0) {$+q$};
  \draw[->,red,very thick] (-4.58,0) -- (-0.58,0);
  \node[red,below] at (-2.00,0) {$v$};

  % 分析磁场与探测器
  \draw[thick,fill=blue!4] (-0.55,-0.20) rectangle (3.20,3.20);
  \foreach \x in {0.05,0.75,1.45,2.15,2.85} {
    \foreach \y in {0.35,1.00,1.65,2.30,2.95} {
      \draw[gray] (\x-0.05,\y-0.05) -- (\x+0.05,\y+0.05);
      \draw[gray] (\x-0.05,\y+0.05) -- (\x+0.05,\y-0.05);
    }
  }
  \node[blue!70!black,font=\small] at (2.42,2.90)
    {分析磁场 $B_2$ 向里};

  \coordinate (S) at (-0.55,0);
  \draw[red,very thick,->] (S) arc (-90:90:0.78);
  \draw[orange!90!black,very thick,->] (S) arc (-90:90:1.28);
  \node[red,font=\small] at (0.30,0.80) {$m_1$};
  \node[orange!90!black,font=\small] at (0.86,1.76) {$m_2>m_1$};

  \draw[very thick] (-0.72,0.20) -- (-0.72,2.85);
  \node[font=\small] at (-0.72,3.04) {位置探测器};
  \fill[red] (-0.55,1.56) circle (2pt);
  \fill[orange!90!black] (-0.55,2.56) circle (2pt);
  \draw[->,red,thin] (-1.75,1.56) -- (-0.67,1.56);
  \node[red,left,font=\small] at (-1.75,1.56) {$x_1=2r_1$};
  \draw[->,orange!90!black,thin] (-1.75,2.56) -- (-0.67,2.56);
  \node[orange!90!black,left,font=\small] at (-1.75,2.56) {$x_2=2r_2$};
\end{tikzpicture}
